\subsection{分配性}

定义5. 设 \(\mathbf{E}_1, \ldots, \mathbf{E}_n\) 和 \(\mathbf{F}\) 为集合， \(u\) 为 \(\mathbf{E}_1 \times \cdots \times \mathbf{E}_n\) 到 \(\mathbf{F}\) 中的一个映射。设 \(i \in [1, n]\) 。假设 \(\mathbf{E}_i\) 和 \(\mathbf{F}\) 被赋予广群结构。称 \(u\) 关于指标变量 \(i\) 是分配的，如果部分映射

\[
x _ {i} \mapsto u (a _ {1}, \dots , a _ {i - 1}, x _ {i}, a _ {i + 1}, \dots , a _ {n})
\]

是 \(\mathbf{E}_i\) 到 \(\mathbf{F}\) 的同态，其中 \(a_j\) 是 \(\mathbf{E}_j\) 中任意固定的元素且 \(j \neq i\) 。

若 \(\top\) 表示 \(\mathbf{E}_i\) 和 \(\mathbf{F}\) 上的内合成律，则 \(u\) 的分配性由下列等式给出：

\[
\begin{array}{r l} & u (a _ {1}, \dots , a _ {i - 1}, x _ {i} \top x _ {i} ^ {\prime}, a _ {i + 1}, a _ {n}) \\ & \quad = u (a _ {1}, \dots , a _ {i - 1}, x _ {i}, a _ {i + 1}, \dots , a _ {n}) \top u (a _ {1}, \dots , a _ {i - 1}, x _ {i} ^ {\prime}, a _ {i + 1}, \dots , a _ {n}) \end{array}\tag{1}
\]

对于 \(i = 1, 2, \ldots, n\) , \(a_1 \in \mathrm{E}_1, \ldots, a_{i-1} \in \mathrm{E}_{i-1}\) , \(x_i \in \mathrm{E}_i\) , \(x_i' \in \mathrm{E}_i\) , \(a_{i+1} \in \mathrm{E}_{i+1}, \ldots, a_n \in \mathrm{E}_n\) .

例。设 E 是乘法记法书写的幺半群（相应地，群）。映射 \((n, x) \mapsto x^{n}\) 把 \(N \times E\) （相应地 \(Z \times E\) ）映入 E ，由等式 \(x^{m+n} = x^{m}x^{n}\) （以加法作为 N 上的律），它关于第一个变量是分配的。若 E 是交换的，则该映射由等式 \((xy)^{n} = x^{n}y^{n}\) 关于第二个变量是分配的。

命题 1. 设 \(\mathbf{E}_1, \mathbf{E}_2, \ldots, \mathbf{E}_n\) 和 \(\mathbf{F}\) 为加法记法书写的交换幺半群，并设 \(u\) 为 \(\mathbf{E}_1 \times \cdots \times \mathbf{E}_n\) 到 \(\mathbf{F}\) 中的一个映射，它关于所有变量都是分配的。对于 \(i = 1, 2, \ldots, n\) ，设 \(\mathbf{L}_i\) 为非空有限集，且 \((x_{i,\lambda})_{\lambda \in \mathbf{L}_i}\) 为 \(\mathbf{E}_i\) 中元素的族。设 \(y_i = \sum_{\lambda \in \mathbf{L}_i} x_{i,\lambda}\) ， \(i = 1, 2, \ldots, n\) 。则

\[
u (y _ {1}, \dots , y _ {n}) = \sum_ {\alpha} u (x _ {1, \alpha_ {1}}, \dots , x _ {n, \alpha_ {n}})\tag{2}
\]

对所有序列求和 \(\alpha = (\alpha_{1},\ldots ,\alpha_{n})\) 的序列 \(\mathbf{L}_1\times \dots \times \mathbf{L}_n\)

我们通过对 \(n\) 进行归纳论证，基础情形 \(n = 1\) 由§1第2节的公式（2）得出。出自同一参考文献。

\[
u (y _ {1}, \dots , y _ {n - 1}, y _ {n}) = \sum_ {\alpha_ {n} \in L _ {n}} u (y _ {1}, \dots , y _ {n - 1}, x _ {n, \alpha_ {n}})\tag{3}
\]

对于 \(y_{n} = \sum_{\alpha_{n} \in \mathbf{L}_{n}} x_{n, \alpha_{n}}\) 且映射 \(z \mapsto u(y_{1}, \ldots, y_{n-1}, z)\) 于 \(\mathbf{E}_{n}\) 到 \(F\) 是一个原群同态。将归纳假设应用于分配映射 \((z_{1}, \ldots, z_{n-1}) \mapsto u(z_{1}, \ldots, z_{n-1}, x_{n, \alpha_{n}})\) 从 \(\mathbf{E}_{1} \times \cdots \times \mathbf{E}_{n-1}\) 到 \(F\) ,

\[
u (y _ {1}, \dots , y _ {n - 1}, x _ {n, \alpha_ {n}}) = \sum_ {\alpha_ {1}, \dots , \alpha_ {n - 1}} u (x _ {1, \alpha_ {1}}, \dots , x _ {n - 1, \alpha_ {n - 1}}, x _ {n, \alpha_ {n}}),\tag{4}
\]

求和是对这些序列进行的 \((\alpha_{1},\ldots,\alpha_{n-1})\) 的序列 \(M=L_{1}\times\cdots\times L_{n-1}\) 的一个子群。现在 \(L_{1}\times\cdots\times L_{n}=M\times L_{n}\) ；记

\[
t _ {\alpha_ {1}, \dots , \alpha_ {n}} = u (x _ {1, \alpha_ {1}}, \dots , x _ {n, \alpha_ {n}}),
\]

我们有

\[
\sum_ {\alpha_ {1}, \dots , \alpha_ {n}} t _ {\alpha_ {1}, \dots , \alpha_ {n}} = \sum_ {\alpha_ {n}} \left(\sum_ {\alpha_ {1}, \dots , \alpha_ {n}} t _ {\alpha_ {1}, \dots , \alpha_ {n - 1}, \alpha_ {n}}\right)\tag{5}
\]

由§1第5节的公式(7)，(2)立即由(3)、(4)和(5)得出。

注记。若 \(u(a_{1},\ldots ,a_{i - 1},0,a_{i + 1},\ldots ,a_{n}) = 0\) 对于 \(i = 1,2,\dots,n\) 和 \(a_{j}\in \mathbf{E}_{j}\) ( \(j\neq i\) )，则公式(2)对有限支撑的族 \((x_{i,\lambda})_{\lambda \in \mathbf{L}_i}\) 仍然成立。

定义5的一个特例是，u为与集合的作用相关联的作用法则的情形。 \(\Omega\) 在岩浆 E 上。如果 u 关于第二个变量是分配的，则也称该作用为 \(\Omega\) 在岩浆 E 上是分配的。换言之：

定义6. 一个作用 \(\alpha \mapsto f_{\alpha}\) 一个集合的 \(\Omega\) 在岩浆 E 上，若对所有 \(\alpha \in \Omega\) ，映射 \(f_{\alpha}\) 是岩浆 E 的一个自同态。

上定义了一个处处有定义的合成律。若 \(\top\) 表示岩浆 E 的运算，而 \(\bot\) 是与 \(\Omega\) 在 E 上的作用相伴的作用运算，后者的分配性则由公式

\[
\alpha \perp (x \top y) = (\alpha \perp x) \top (\alpha \perp y) \quad (\text { for   } \alpha \in \Omega \text {   and   } x, y \in E).\tag{6}
\]

表达。借用语言上的习惯，也称运算 \(\perp\) 关于运算 \(\top\).

是分配的（或右分配的）。§1 第 2 号中的公式 (2) 表明，此时，对 E 中元素的每个有序序列 \((x_{\lambda})_{\lambda \in L}\) 以及每个 \(\alpha \in \Omega\) ,

\[
\alpha \perp \left(\mathsf {T} _ {\lambda \in L} x _ {\lambda}\right) = \mathsf {T} _ {\lambda \in L} (\alpha \perp x _ {\lambda}).\tag{7}
\]

如果一个作用 \(\alpha\mapsto f_{\alpha}\) 是可分配的，并且 E 上的等价关系 R 与 E 的合成律及该作用相容 \(\alpha\mapsto f_{\alpha}\) ，则 E/R 上的商作用是可分配的。

当 E 上的律用乘法书写时，我们常使用指数记号 \(x^{\alpha}\) 来表示关于该乘法可分配的作用律，因此可分配性由恒等式 \((xy)^{\alpha} = x^{\alpha}y^{\alpha}\) 表达。若 E 上的律用加法书写，我们常使用左（相应地，右）乘法记号 \(\alpha.x\) （相应地 \(x.\alpha\) ）来表示关于该加法可分配的作用律，其可分配性由恒等式

\[
\alpha (x + y) = \alpha x + \alpha y \quad (\text { resp. } (x + y) \alpha = x \alpha + y \alpha).
\]

表达。我们也可以考虑 \(\Omega\) 具有一个内律，记为 \(\overline{\top}\) ，且作用律关于第一个变量可分配，这意味着

\[
(\alpha \overline{\top} \beta) \perp x = (\alpha \perp x) \top (\beta \perp x)\tag{8}
\]

为所有 \(\alpha, \beta \in \Omega\) 和 \(x \in \mathbf{E}\) 。于是，由 §1 第 2 号的公式 (2)

\[
\left(\overline {{\mathsf {T}}} _ {\lambda \in L} \alpha_ {\lambda}\right) \perp x = \mathsf {T} _ {\lambda \in L} (\alpha_ {\lambda} \perp x)\tag{9}
\]

对每一个有序序列 \((\alpha_{\lambda})_{\lambda \in L}\) 的 \(\Omega\) 以及所有 \(x\in E\) .

